\documentclass[10pt,a4paper]{amsart}
\usepackage[margin=19mm]{geometry}
\usepackage{amsmath,amssymb,mathtools}
\usepackage[scr=rsfs,cal=euler]{mathalfa}
\usepackage{xfrac}
\usepackage[unicode,hidelinks,bookmarks=false]{hyperref}
\newcommand{\ie}{i.e.\ }
\newcommand{\cf}{cf.\ }
\newcommand{\resp}{respectively\ }
\newcommand{\Subsection}{\subsection}
\providecommand{\nobreakdash}{\nobreak\hbox{-}}
\setlength{\emergencystretch}{2em}
\multlinegap0pt
\usepackage{amssymb}
\usepackage{mathtools}
\usepackage{xcolor}
\newtheorem{lemma}{Lemma}
\usepackage{cleveref}
\crefname{lemma}{Lemma}{Lemmas}
\newcommand{\T}{\mathbb T}
\newcommand{\Z}{\mathbb Z}
\newcommand{\one}{\mathbf 1}
\title{Obstructions to coloring arithmetic graphs}
\author{Lujia Wang, Ruihua Wang}
\date{Source: arXiv:2609.15081v1 (CC BY 4.0)}
\begin{document}
\maketitle

\noindent\emph{Source and licence.} Obstructions to coloring arithmetic graphs. Version arXiv:2609.15081v1 (CC BY 4.0), distributed by arXiv under the Creative Commons Attribution 4.0 International licence, http://creativecommons.org/licenses/by/4.0/, as recorded in the arXiv licence metadata for this exact version and shown on the abstract page https://arxiv.org/abs/2609.15081v1. Full licence notice: \texttt{licenses/arxiv-2609.15081-CC-BY-4.0.txt}.\par\smallskip

\noindent\textit{Editorial scope.} Counted unit: the article's lemma lem:205-phases (source0.tex lines 563--571). The article's complete original proof of it is delivered with it (source0.tex lines 573--813, one \texttt{proof} environment of 241 lines). No mathematics is written, repaired or re-derived: the statement and the proof are the article's own lines, in the article's own notation, and no retained line is substituted or re-flowed.  Retained uncounted as the source-local context the proof needs: lines 353--413.  Nothing was omitted from any retained span, so the package carries no omission record.  No retained line contains a citation, so the wrapper carries no bibliography and no bibliography entry.  No retained line contains a figure environment, an image-inclusion command or drawing code, so no image is delivered and the package declares no asset dependency.  Editorial formatting only, in addition: an AMS \texttt{amsart} preamble that restates the article's own declarations and macros used by the retained lines, namely the environments \texttt{lemma} on separate counters, together with the standard packages those lines need.  The retained environments print in this document as Lemma 1 in order of appearance, whereas the article prints the same items with the numbering of its own full document.  The complete native source of the article as distributed in the arXiv e-print for this exact version is bundled unchanged as \texttt{sources/210387/source0.tex}.
\begin{lemma}[Finite arithmetic checks; computer-assisted]\label{lem:205-computations}
The following statements hold.
\begin{enumerate}
\renewcommand{\theenumi}{\roman{enumi}}
\renewcommand{\labelenumi}{\textup{(\theenumi)}}
\item There is no homomorphism $(u,v):\Z^d\to\Z_{125}\times\Z_{41}$
such that, for every $b\in\Z_{41}$, the labels with $v(t(m))=b$ have
five distinct $u$-values forming one coset of $25\Z_{125}$.
Equivalently, each fiber contains exactly five labels, its $u$-values
agree modulo $25$, and no two agree modulo $125$.

\item For each $s\in\{0,1,2,3\}$, there is no choice of weights in
$\Z\times\Z_{41}$ for which $W([205])$ is supported on
$\{0,1,2,3\}\times\Z_{41}$, with
\[
 \mu_W(h,b)=1+\one_{\{h=s\}},
 \qquad (h,b)\in\{0,1,2,3\}\times\Z_{41},
\]
and weights
\[
 w_2=(0,1),\quad w_3=(0,b_3),\quad w_5=(1,0),\quad
 w_r\in\{0,1\}\times\Z_{41}\quad(r=7,11,13).
\]
For each $s\in\{4,5\}$, there is likewise no choice of weights in
$\Z_6\times\Z_{41}$ with
\[
 \mu_W(h,b)=\one_{\{h\ne s\}},\qquad
 w_2=(0,1),\quad w_3=(0,b_3),\quad w_5=(1,b_5).
\]
Here $b_3,b_5\in\Z_{41}$ are free, as are the other prime weights
unless explicitly restricted.

\item If $v:\Z^6\to\Z_{41}$ satisfies $v(e_2)=1$ and every
fiber of $v|_{T_0}$ has at most five labels, then $D_0(v)$ is
$5$-saturated in $\Z^6$:
\[
 5x\in D_0(v)\ \Longrightarrow\ x\in D_0(v)
 \qquad(x\in\Z^6).
\]
Equivalently, every nonzero invariant factor of $D_0(v)\le\Z^6$
is relatively prime to $5$.

\item Consider homomorphisms $L:\Z^6\to\Z_{205}$ whose
restrictions to $T_0$ are injective. Reducing $L$ modulo $41$ and
multiplying by a nonzero scalar to obtain $v(e_2)=1$ gives precisely
$51\,582$ distinct core maps $v$. For $51\,451$ of these,
$\operatorname{rank}D_0(v)=6$; the other $131$ have rank five.
For each rank-five map there is a primitive integer linear functional
$H:\Z^6\to\Z$ annihilating $D_0(v)$, with $H_r=H(e_r)$, such that
\[
 H_2\not\equiv0\pmod{41},\qquad H_r\equiv H_2v(e_r)\pmod{41},
\]
and
\[
 \max_{t\in T_0}H(t)-\min_{t\in T_0}H(t)\le58,
 \qquad
 \max_{1\le c\le12}H(t(c))-\min_{1\le c\le12}H(t(c))\le29.
\]
Here primitive means $\gcd(H_2,H_3,H_5,H_7,H_{11},H_{13})=1$.
\end{enumerate}
\end{lemma}

\begin{lemma}[Character values at small primes]\label{lem:205-phases}
Suppose the multiset $z(T)$ is a union of five complete $41$-gons.
Then
\[
 \operatorname{ord}(\alpha_2),\operatorname{ord}(\alpha_3),
 \operatorname{ord}(\alpha_5)\le4,
 \qquad z_2^{12}\in\mu_{41}\setminus\{1\}.
\]
\end{lemma}

\begin{proof}
We will repeatedly use the following consequence of
\cref{lem:205-computations}(i).
Suppose integer prime weights $k_r$ define the additive height
$K(x)=\sum_r k_rx_r$, whose values on $T$ lie in five consecutive
integers, and suppose these five heights distinguish the five values in $E$.
If a character $\beta:\Z^d\to\mu_{41}$ is obtained from $z$ by a
constant rotation on each height, then every height contains
each value of $\mu_{41}$ exactly once. Fix
$\zeta_{41}=e^{2\pi i/41}$ and write
$\beta(x)=\zeta_{41}^{b(x)}$, with $b:\Z^d\to\Z_{41}$ additive.
The homomorphism
\begin{equation}\label{eq:205-height-logarithm}
 x\longmapsto \bigl(K(x)\bmod5,b(x)\bigr)
       \in\Z_5\times\Z_{41}\cong\Z_{205}
\end{equation}
is then bijective on $T$: different heights have different first
coordinates, and within each height the second coordinates run through
$\Z_{41}$ once. Its restriction to the exponent labels is therefore
a logarithm of length $205$, contrary to
\cref{lem:205-computations}(i). The requirement that $K$ be an
\emph{additive integer} height matters: exponents chosen modulo the
order of a character value need not lift to such a height.

\smallskip\noindent
\emph{The values at $2$ and $3$.}
The eight labels $2^j$, $0\le j\le7$, show that $\alpha_2$ has finite order
at most five: an element of larger or infinite order would already
give six distinct values among its first six powers.
Suppose $u=\alpha_3$ has order $N>5$, including $N=\infty$.
The five labels $1,3,9,27,81$ then exhaust the set $E$:
\[
 E=\{1,u,u^2,u^3,u^4\}.
\]
Put $U=\langle\alpha_2\rangle$. Its order is at most five, and
the labels $2^j,3\cdot2^j,9\cdot2^j$, for $0\le j<|U|$, all
belong to $[205]$. Hence $U,uU,u^2U\subset E$.
This forces $U=\{1\}$. Indeed, $U\subset\langle u\rangle$ is
finite, so the assertion is immediate when $N=\infty$. For finite
$N\ge9$, a nontrivial subgroup of order $s$ contains
$u^{N-N/s}$, whose exponent is at least five and less than $N$.
For $N=6$, the possible nontrivial subgroups have order two or three:
$u^2U$ contains $u^5$ in the former case, and $uU$ contains $u^5$
in the latter. For $N=7$ there is no nontrivial subgroup of order at
most five. For $N=8$, the subgroup of order four already contains
$u^6$, while the subgroup of order two has $u^5\in uU$.
All possibilities contradict the displayed description of $E$.
Thus $\alpha_2=1$.

For each prime $r\le205$, choose the unique $k_r\in\{0,1,2,3,4\}$
such that $\alpha_r=u^{k_r}$. In particular $k_2=0$ and $k_3=1$.
For $r=5,7,11,13$, the labels $3r$ and $9r$ first force $k_r\le2$:
the choices three and four would produce the forbidden value $u^5$.
The labels $75=3\cdot5^2$ and $147=3\cdot7^2$ further exclude
$k_5=k_7=2$. Consequently
\[
 k_2=0,\quad k_3=1,\quad k_5,k_7\le1,\quad
 k_{11},k_{13}\le2.
\]
Define $K(x)=\sum_r k_rx_r$ using these integer weights.
For a core integer $m$, each core prime satisfies $r\ge3^{k_r}$,
so $m\ge3^{K(t(m))}$. Since $205<3^5$, this gives
$0\le K(t(m))\le4$ on the core.
For a tail $m=rc$, with $r\ge17$ prime and $c\le12$, the preceding
weights give $K(t(c))\le2$. If this cofactor height is one, then
$c\ge3$ and the available label $3r$ forces $k_r\le3$.
If it is two, then $c\ge9$ and the labels $3r,9r$ force $k_r\le2$.
These deductions use only exponents up to five, so no reduction modulo
$N$ occurs, even when $N=6$. Thus $K(t(m))\le4$ on every tail too.

All five heights occur at powers of three. Set
$\beta_r=z_r/z_3^{k_r}$ and extend multiplicatively to a character.
Since $\beta_r^{41}=\alpha_r/u^{k_r}=1$, this character takes values
in $\mu_{41}$. At height $j$ it divides every tile value by $z_3^j$,
so \eqref{eq:205-height-logarithm} gives the forbidden logarithm.
We conclude that $\operatorname{ord}(\alpha_3)\le5$.

If either $\alpha_2$ or $\alpha_3$ has order five, its observed powers
give $E=\mu_5$. Then $z(T)\subset\mu_{205}$, and its five gons are
the five distinct $\mu_{41}$-cosets of $\mu_{205}$, each used once.
The character $z$ itself therefore gives a logarithm to
$\mu_{205}\cong\Z_{205}$. This excludes order five at both primes.

The resulting bound $\operatorname{ord}(\alpha_2)\le4$ gives
$z_2^{12}\in\mu_{41}$. Suppose it were one. All eight values
$z(t(2^j))=z_2^j$ would then belong to $\mu_{12}$. Two such values in
the same $\mu_{41}$-coset must be equal, because
$\mu_{12}\cap\mu_{41}=\{1\}$. But one complete gon contains each
of its values only once. Hence each of the eight labels would need a
different gon, even if some values repeat. Five gons cannot suffice,
and $z_2^{12}\ne1$ follows.

\smallskip\noindent
\emph{The subgroup $\langle\alpha_2,\alpha_3\rangle$.}
Let $H=\langle\alpha_2,\alpha_3\rangle$. A finite subgroup of
$\T$ is cyclic, so its order is the least common multiple of the
two generator orders. Since both orders are at most four, the only
possibilities exceeding four are six and twelve. In those cases the
generator orders are respectively $2,3$ or $3,4$, in either order.
All their distinct products occur among $2^a3^b\le108$, contradicting
$|E|\le5$. Thus $|H|\le4$.

If $|H|=s\in\{3,4\}$, one of $\alpha_2,\alpha_3$ generates $H$.
Writing that prime as $r$, the labels $r^j,5r^j$ for $0\le j<s$
are at most $5\cdot3^3=135$ and show that
$H\cup\alpha_5H\subset E$. Distinct cosets would contribute
$2s>5$ values, so $\alpha_5\in H$.
If $|H|=2$, choose $r\in\{2,3\}$ with $\alpha_r\ne1$.
The six labels $r^i5^j$, $0\le i<2$, $0\le j<3$, are at most
$75$ and contain the three cosets $H,\alpha_5H,\alpha_5^2H$.
They cannot be distinct. The order of $\alpha_5H$ in
$\langle H,\alpha_5\rangle/H$ is therefore at most two, and the
whole group has order at most four. We have proved
\begin{equation}\label{eq:205-small-phase-subgroup}
 H\ne\{1\}\quad\Longrightarrow\quad
 |\langle\alpha_2,\alpha_3,\alpha_5\rangle|\le4.
\end{equation}
In particular, any remaining case with
$\operatorname{ord}(\alpha_5)>4$ has $\alpha_2=\alpha_3=1$.

\smallskip\noindent
\emph{The four initial values at $5$.}
Put $u=\alpha_5$ and assume its order is at least five or infinite.
The labels $1,5,25,125$ give the four distinct values
$I=\{1,u,u^2,u^3\}\subset E$. There are only two possibilities:
$E=I$, or $E=I\cup\{v\}$ for a single $v\notin I$.

First consider $E=I$. Write $\alpha_r=u^{k_r}$ with
$0\le k_r\le3$. For $r=7,11,13$, the labels $5r$ and $r^2$
exclude respectively $k_r=3$ and $k_r=2$, since $u^4\notin E$.
Thus the core weights are
\[
 k_2=k_3=0,\qquad k_5=1,\qquad
 k_7,k_{11},k_{13}\in\{0,1\}.
\]
For a core label, height at least four would require an integer at
least $5^4>205$. For a tail $rc$, the cofactor height is at most one,
so the unreduced height $k_r+K(t(c))$ is at most four. Its value
cannot be four because $u^4\notin E$. Thus the additive integer height
$K$ takes values only in $\{0,1,2,3\}$ on $T$. Each height corresponds
to at least one complete gon; as there are five gons in total, one height contains two
and the other three contain one each.

Set $\beta_r=z_r/z_5^{k_r}$ and write
$\beta_r=\zeta_{41}^{b_r}$. Here $b_5=0$. Also $\beta_2=z_2$ is a
nontrivial $41$st root: $\alpha_2=1$ and $z_2^{12}\ne1$.
Multiplying every $b_r$ by $b_2^{-1}$ normalizes $b_2=1$ and preserves
all cell multiplicities. The resulting arithmetic image in
$\{0,1,2,3\}\times\Z_{41}$ is exactly the doubled-height model
excluded by \cref{lem:205-computations}(ii). Thus $E\ne I$.

\smallskip\noindent
\emph{An additional value.}
Suppose now $E=I\cup\{v\}$. If
$v\notin\{u^4,u^{-1}\}$, the equations from $r,5r,r^2$ for
$r=7,11,13$ imply $\alpha_r\in\{1,u\}$.
Indeed, $u^2$ would require the missing square $u^4$, and $u^3$
would require the missing product $u^4$. A prime with value $v$ and
$5r\le205$ would require $uv\in E$; since $v\notin I$, this is
possible only when $uv=1$, or $v=u^{-1}$, which is excluded.
Consequently every prime with value $v$ exceeds $41$.

Give such exceptional primes height four. Give every other prime
the height $k_r\in\{0,1,2,3\}$ determined by $\alpha_r=u^{k_r}$.
A label containing an exceptional prime contains it once, has cofactor
at most four, and hence has height four and value $v$.
For all other labels, the same core and tail estimates as above bound
the unreduced height by four; height four is excluded by $u^4\notin E$.
Thus these labels have heights $0,1,2,3$. The five heights correspond
bijectively to the five values in $E$.

Choose $\theta\in\T$ with $\theta^{41}=v$ and put
\[
 \beta_r=
 \begin{cases}
 z_r/\theta,&\alpha_r=v,\\
 z_r/z_5^{k_r},&\alpha_r\ne v.
 \end{cases}
\]
These prime values lie in $\mu_{41}$ and define one character
$\beta$. On heights $0,1,2,3$ its denominator is $z_5^K$; on height
four its denominator is $\theta$, because the exceptional prime occurs
once and its cofactor has height zero. Thus each height still contains
a complete $41$-gon after division. The logarithm
\eqref{eq:205-height-logarithm} is again a contradiction.
This argument applies to every order at least five, including infinite
order, whenever the additional value is neither $u^4$ nor $u^{-1}$.

It remains to consider $v=u^4$ and $v=u^{-1}$. First suppose
$\operatorname{ord}(u)\ge7$, or $u$ has infinite order.
If $v=u^4$, use prime heights in $[0,4]$.
The labels $5r,r^2$ show that the core primes $r=7,11,13$ have
heights at most two: height three would require both $u^4$ and $u^6$,
and height four would require $u^5$. A core label of height at least
five would therefore have size at least
\[
 \min\{5^5,5^3\cdot7,5\cdot7^2,7^3\}=245>205.
\]
A tail has cofactor height at most two and hence unreduced height at
most six. These seven powers are distinct, so membership in $E$ forces
the height to be at most four.
If $v=u^{-1}$, use prime heights in $[-1,3]$. A core prime cannot
have height $-1$, since its square would have the absent value $u^{-2}$;
heights two and three would require the absent value $u^4$ at $r^2$
or $5r$. Thus $k_7,k_{11},k_{13}\in\{0,1\}$.
Core heights are in $[0,3]$, and a tail has unreduced height in
$[-1,4]$. These powers are distinct; since $u^4\notin E$, its height
lies in $[-1,3]$. In both cases the five consecutive heights
are attained and correspond to the five values in $E$. Dividing by $z_5^K$
produces a character in $\mu_{41}$ and the forbidden logarithm
\eqref{eq:205-height-logarithm}. Together with the preceding cases,
this proves $\operatorname{ord}(\alpha_5)\le6$.

For order five, an additional value in $\mu_5$ makes $E=\mu_5$,
which directly gives the logarithm already excluded above. An additional
value outside $\mu_5$ is covered by the exceptional-prime construction,
and $E=I$ was excluded by the doubled-height computation.
For order six, the same arguments exclude $E=I$ and an additional
value outside $\mu_6$. The two remaining possibilities are
\[
 E=\mu_6\setminus\{u^4\}
 \quad\hbox{or}\quad E=\mu_6\setminus\{u^5\}.
\]
Use a coordinate in $\Z_6$. Since $\gcd(41,6)=1$, choose
$\theta\in\mu_6$ with
$\theta^{41}=u$. There is an additive map
$k:\Z^d\to\Z_6$ such that $\alpha(x)=u^{k(x)}$, and
$\beta(x)=z(x)/\theta^{k(x)}$ is a well-defined character into
$\mu_{41}$. Writing $\beta=\zeta_{41}^b$ gives an arithmetic image
in $\Z_6\times\Z_{41}$. Each of the five permitted heights contains
one complete gon, so every cell outside the one omitted height is
filled exactly once. As before, normalize $b_2=1$; the prime weights
have the form
\[
 w_2=(0,1),\qquad w_3=(0,b_3),\qquad w_5=(1,b_5).
\]
These are precisely the two six-height models excluded by
\cref{lem:205-computations}(ii). Here $b_5$ is unrestricted, since
$\theta$ must lie in $\mu_6$ whereas $z_5$ need not.
Orders five and six are excluded, completing the proof.
\end{proof}

\end{document}
